\chapter{Pollen allergy}
\index{Pollen allergy}

\medskip
\noindent\textit{Educational information; seek professional advice for individual care.}

\section*{What it is}
Pollen allergy, often called hay fever or seasonal allergic rhinitis, occurs when the immune system reacts to pollen from trees, grasses, or weeds. It is different from a cold, although symptoms can look similar.

\section*{Symptoms}
Sneezing, an itchy or blocked nose, watery or itchy eyes, an itchy throat or palate, cough, reduced smell, headache, and tiredness are common. Symptoms often vary with the local pollen season and weather, but can occur for longer depending on exposure.

\section*{Management}
A pharmacist or clinician can advise about non-drowsy antihistamines, antihistamine eye drops, or steroid nasal sprays. Use nasal sprays with the correct technique and follow age, pregnancy, and medicine-safety advice. Reduce exposure by keeping windows closed when counts are high, showering after outdoor activity, and changing clothes; avoid measures that worsen asthma or wellbeing.

\section*{When to seek medical care}
Arrange review if symptoms are severe, persistent, affecting sleep or daily activities, or not controlled with suitable pharmacy treatment. Seek urgent help for severe breathing difficulty, swelling of the lips or tongue, fainting, or a rapidly worsening allergic reaction.

\section*{Sources}
\begin{itemize}
\item NHS hay fever guidance: \url{https://www.nhs.uk/conditions/hay-fever/}
\item NHS allergic rhinitis guidance: \url{https://www.nhs.uk/conditions/allergic-rhinitis/}
\end{itemize}
